\documentclass[11pt]{article}
\usepackage[margin=0.75in]{geometry}
\usepackage{fontspec}
\usepackage{setspace}
\usepackage{hyperref}
\hypersetup{colorlinks=true, urlcolor=blue, linkcolor=black}
\pagestyle{empty}
\setlength{\parskip}{0.48em}
\setlength{\parindent}{0pt}
\setstretch{0.96}
\setmainfont{Times New Roman}

\begin{document}

\begin{flushleft}
Ke Cheng HSU (Clio) \\
Department of Economics, The Hong Kong University of Science and Technology \\
Clear Water Bay, Kowloon, Hong Kong $\cdot$ +852 6840 9753 \\
\href{mailto:khsu@connect.ust.hk}{khsu@connect.ust.hk} $\cdot$ \href{https://cliohsu1997.github.io/}{cliohsu1997.github.io}
\end{flushleft}

7 November 2026

\vspace{0.2em}
Faculty Search Committee \\
Department of Economics, National Taipei University \\
\href{mailto:econ_job@mail.ntpu.edu.tw}{econ\_job@mail.ntpu.edu.tw}

\vspace{0.2em}
Dear Members of the Search Committee,

I am writing to apply for the full-time Assistant Professor position in the Department of Economics at National Taipei University. I am a Ph.D.\ candidate in Economics at the Hong Kong University of Science and Technology. I expect to graduate in June 2027, before the 31 July 2027 degree deadline, and I can begin on 1 August 2027. My fields are Empirical Industrial Organization and Digital Economics. I can teach in both Chinese and English.

My job market paper, ``Designing Return Policy in E-Commerce: A Trade-off between Purchase and Return Efficiency,'' studies whether a return policy encourages efficient returns and purchases. A stricter policy curbs inefficient returns after delivery, but the lower return rate makes a higher product price more attractive for the seller and can exclude efficient purchases. Using transaction data from a Taobao apparel retailer, I estimate a structural model of purchase, return, and whether to buy return-shipping insurance under monopoly pricing across six categories. Automatic coverage of return shipping performs better when the margin, net scrap value, and uncertainty are high, because it expands efficient purchases. Optional coverage performs better when those attributes are low, because it curbs excessive returns. A menu that pairs a discounted price with a medium restocking fee balances the two margins and performs better than both insurance policies for most categories.

A second working paper, ``Product Design with Heterogeneous Search Costs'' (with Pak Hung Au), studies how firms choose product design when consumers differ in search costs. Search costs shape which products consumers search. Consumers with low search costs seek niche products, and consumers with high search costs seek standardized products. A fall in search costs can raise prices when this shift in which product consumers search outweighs the effect that encourages more search. I also have two ongoing projects. With Kohei Kawaguchi and Kei Kawai, I use Japanese bookstore data to study how resale price maintenance encourages retailers to hold inventory in low-demand states. However, unsold copies cannot be cleared by cutting the price, so they are returned and discarded at an environmental cost. With Ce Bian and Liang Zhao, I study pass-through when products can be returned, and show that a tax can lower the equilibrium price and raise welfare by reducing the return rate.

At HKUST I was a teaching assistant from 2020 to 2025 for undergraduate macroeconomics, the economics of uncertainty and information, and Ph.D.\ microeconomic theory, with an average evaluation of 4.72/5. In 2026 I designed and lectured the mathematics camp for the MSc in Economics, receiving 4.53/5. I am prepared to teach microeconomic theory, industrial organization, econometrics, and applied courses that use large transaction data, from the undergraduate level through the graduate level.

I attach my CV, job market paper, a second research paper, my transcript, proof of my expected graduation date, evidence of teaching, and the two faculty forms. The official transcript will be ready about two months before I graduate in June 2027. The courses and grades will largely be the same as on the transcript I submit now, and I will send the official transcript as soon as I receive it. My referees will send their letters separately. You can reach me at khsu@connect.ust.hk or +852 6840 9753.

\vspace{0.45em}
Sincerely,

\vspace{0.45em}
Ke Cheng HSU \\
Ph.D.\ Candidate, Department of Economics, HKUST

\end{document}
